% !PRINTED 242
\weberpage{242}
%<<CONT
gehört. Da mit den Derivirten $f'(\xi+i\eta), f''(\xi+i\eta)\ldots$ zugleich die conjugirten Grössen $f'(\xi-i\eta), f''(\xi-i\eta)\ldots$ verschwinden, so folgt, dass \gesperrt{conjugirt imaginäre Wurzeln denselben Grad der Vielfachheit haben.}

Daraus folgt, dass die imaginären Wurzeln immer in gerader Zahl vorkommen, und dass eine Gleichung ungeraden Grades immer mindestens eine reelle Wurzel haben muss.

Unser nächstes Ziel wird das sein, die Zahl der reellen Wurzeln, ohne die Gleichung aufzulösen, direct aus den Werthen \gesperrt{gewisser rationaler Functionen der Co\"efficienten der Gleichung} zu bestimmen.

Eine sehr wichtige Rolle spielt hierbei die Discriminante, auf deren Bedeutung für unsere Frage wir zunächst eingehen müssen.

Wir haben in \S.\,46 die Discriminante erklärt als das Product aus den Quadraten der sämmtlichen Wurzeldifferenzen
\[
(x_1 - x_2)^2,
\]
noch multiplicirt mit $a_0^{2n-2}$, wenn $a_0$ der Co\"efficient der höchsten Potenz der Unbekannten in der gegebenen Gleichung ist, und wir haben dort gezeigt, wie die Discriminante als rationale Function der Co\"efficienten berechnet werden kann. Der Factor $a_0^{2n-2}$ ist immer positiv und könnte hier ohne wesentliche Beschränkung der Allgemeinheit auch gleich 1 vorausgesetzt werden. Wir wollen, wie schon früher, für die Discriminante das Zeichen $D$ gebrauchen.

\gesperrt{Die Discriminante verschwindet dann und nur dann, wenn unter den Wurzeln zwei gleiche vorkommen.}

Nehmen wir an, dass $f(x)$ durch Absonderung des grössten gemeinschaftlichen Theilers von $f(x)$ und $f'(x)$, was durch rationale Rechnung geschieht, von mehrfachen Factoren befreit sei, so wird also $D$ nicht verschwinden.

Sind $x_1$ und $x_2$ reell, so ist $(x_1 - x_2)^2$ positiv. Ist $x_1$ reell und $x_2$ imaginär, so giebt es eine zu $x_2$ conjugirte Wurzel $x_2'$, und das Product
\[
(x_1 - x_2)(x_1 - x_2')
\]
ist als Product zweier conjugirt imaginärer Grössen positiv, also auch sein Quadrat.
